\documentclass[11pt]{article}
\usepackage{amsmath,amsthm,amssymb}
\usepackage[colorlinks,citecolor=blue,urlcolor=blue]{hyperref}
\usepackage{tcolorbox}

\newtheorem{theorem}{Theorem}
\newtheorem{corollary}{Corollary}
\newtheorem{lemma}{Lemma}
\newtheorem{remark}{Remark}

\title{A Dissipative Flow for the Critical Strip\\
       and an Exponential Convergence Lemma}
\author{Brayden Sanders\\7Site LLC\\
\small\href{https://doi.org/10.5281/zenodo.18852047}%
{Zenodo DOI 10.5281/zenodo.18852047}}
\date{March 2026}

\begin{document}
\maketitle

\footnotetext{Companion paper: ``Survivor Lines in Finite Affine Planes
$\mathrm{AG}(2,p)$,'' arXiv:2026.XXXXX [math.CO].}

\begin{abstract}
We provide a dynamical-systems reinterpretation of the analytic behaviour
of the Riemann zeta function.
Define the scalar flow $\dot\sigma=-(\sigma-\tfrac12)|\zeta(\sigma+it_0)|^2$
on the critical strip.
This flow has two kinds of fixed points: the critical line $\sigma=\tfrac12$
and the non-trivial zeros of $\zeta$.
We prove unconditionally that in the Korobov--Vinogradov zero-free strip the
flow converges exponentially to $\sigma=\tfrac12$ with the explicit rate
$m_{\mathrm{KV}}(t_0)=\bigl(C_{\mathrm{KV}}(\log t_0)^{2/3}\bigr)^{-2}$.
Consequently,
\[
  \text{RH}\;\Longleftrightarrow\;
  m(t_0):=\min_{0\le\sigma\le1}|\zeta(\sigma+it_0)|^2>0
  \;\text{ for every }t_0>0,
\]
turning the Riemann Hypothesis into a quantitative positivity problem
for a one-dimensional dissipative flow.
\end{abstract}

\section{Introduction}

The Riemann Hypothesis (RH) asserts that all non-trivial zeros of the
Riemann zeta function $\zeta(s)$ satisfy $\operatorname{Re}(s)=\tfrac12$.
We offer a new geometric reformulation: the zeros are the \emph{fixed
points} of a dissipative flow on the critical strip $0<\sigma<1$, and
RH is equivalent to these fixed points all lying on the line
$\sigma=\tfrac12$.

The flow is motivated by the finite-dimensional TIG (Trinity Infinity
Geometry) absorption algebra~\cite{tig}, in which a two-step convergence
theorem plays the same structural role as the Halving Lemma below.
The present note abstracts that finite result to the analytic setting.

\medskip\noindent\textbf{Motivating analogy.}
Both systems share the same dynamical architecture.
In TIG, a fixed-width \emph{shell} of half-width $W=3/50$ around the
mid-plane $3.5$ defines the zone where the algebra provably enforces
return in at most two steps; anything straying beyond $\pm 3W$ is pulled
back to the attractor immediately.
In the $\zeta$-flow, the Korobov--Vinogradov zero-free collar plays the
same role: its half-width $\Delta\sigma_{\mathrm{KV}}(t)\sim
c(\log t)^{-2/3}$ is the zone where Theorem~\ref{thm:kv} gives
unconditional exponential return to $\sigma=\tfrac12$.
The two bands are the same \emph{conceptual} object—an intrinsic safety
zone around the attractor axis—but differ in character: TIG's is discrete
and constant-width; the $\zeta$ collar is analytic and height-dependent,
contracting as $t\to\infty$.
Numerically they coincide only at small heights; thereafter the analytic
window narrows while $W$ stays fixed.
Just as a detector's response band limits which excitations are visible
without secondary scattering, the $\zeta$-flow shows that only starting
points within a shrinking collar around $\sigma=\tfrac12$ maintain a slow
drift; everything outside is driven swiftly back.
Completing RH means pushing that collar all the way inward to
$\sigma=\tfrac12$ itself—exactly what TIG already achieves for its own
attractor.
In the $\zeta$-world the ``shell radius'' is the Korobov--Vinogradov
distance $\Delta\sigma_{\mathrm{KV}}(t_0)\sim c(\log t_0)^{-2/3}$,
which shrinks with height.
The TIG wobble quantum $W=3/50$ is a fixed toy scale in a 9-point
universe; in $\zeta$ the corresponding inner collar shrinks as
$(\log t)^{-2/3}$.
The numerical near-coincidence $W\approx\Delta\sigma_{\mathrm{KV}}$
holds only near $t_0\approx10$ and is a mnemonic, not a prediction.
What is structurally solid is the architecture: a nucleus ($\sigma=\tfrac12$),
a proven inner collar (Korobov--Vinogradov), and an open innermost gap
whose positivity is equivalent to the Riemann Hypothesis.
The analytic question below asks whether that inner gap can be closed—
exactly as the TIG algebra forces its residuals to appear only in their
anchor columns.

\medskip
\begin{tcolorbox}[colback=blue!5,colframe=black!70,
  title={The analytic open problem (one line)}]
Does there exist an absolute constant $c>0$ such that for every $t_0>0$
with no zero of $\zeta$ on $\{{\sigma+it_0:0\le\sigma\le1}\}$,
\begin{equation}\label{eq:open}
  \min_{0\le\sigma\le1}|\zeta(\sigma+it_0)|
  \;\ge\;
  e^{-c(\log t_0)^{2/3}(\log\log t_0)^{1/3}}\,?
\end{equation}
If yes: the Halving Lemma gives exponential convergence for every
zero-free vertical and RH follows.
If no: the explicit failure height $t_0$ locates precisely where the
$\zeta$-flow picture must be refined.
\end{tcolorbox}
\medskip

\section{Statements of Results}

For fixed $t_0>0$, define the scalar ODE on $[0,1]$:
\begin{equation}\label{eq:flow}
  \dot\sigma \;=\; -\bigl(\sigma-\tfrac12\bigr)\,|\zeta(\sigma+it_0)|^2,
  \qquad \sigma(0)=\sigma_0\in[0,1].
\end{equation}

\begin{remark}[Well-posedness]\label{rem:wp}
Since $t_0>0$ the segment $\{\sigma+it_0:0\le\sigma\le1\}$ avoids the
pole of $\zeta$ at $s=1$, so $\zeta$ is analytic there.
The right-hand side of~\eqref{eq:flow} is locally Lipschitz, giving a
unique global solution via Picard--Lindel\"of.
At $\sigma=0$ the drift is $+|\zeta(it_0)|^2/2>0$; at $\sigma=1$ it is
$-|\zeta(1+it_0)|^2/2<0$; hence $[0,1]$ is positively invariant.
\end{remark}

\noindent\textbf{Unconditional theorem (Korobov--Vinogradov strip).}

\begin{theorem}[Explicit convergence in zero-free strip]\label{thm:kv}
Let $t_0\ge3$ and $\sigma_{\mathrm{KV}}(t_0)=1-c(\log t_0)^{-2/3}
(\log\log t_0)^{-1/3}$ denote the Korobov--Vinogradov zero-free boundary.
For any $\sigma_0\in[\sigma_{\mathrm{KV}}(t_0),1]$,
\[
  |\sigma(t)-\tfrac12|
  \;\le\;
  |\sigma_0-\tfrac12|\,
  e^{-m_{\mathrm{KV}}(t_0)\,t},
  \qquad
  m_{\mathrm{KV}}(t_0)
  :=\frac{1}{\bigl(C_{\mathrm{KV}}\,(\log t_0)^{2/3}\bigr)^2}>0,
\]
where $C_{\mathrm{KV}}$ is Ford's explicit constant~\cite[Thm.\,2]{ford}.
\end{theorem}

\noindent\textbf{Equivalence corollary.}

\begin{tcolorbox}[colback=gray!10,colframe=black,title=Equivalence]
\begin{corollary}[RH $\Leftrightarrow$ uniform positivity]\label{cor:equiv}
The Riemann Hypothesis holds if and only if
$m(t_0):=\min_{0\le\sigma\le1}|\zeta(\sigma+it_0)|^2>0$ for every $t_0>0$.
\end{corollary}
\end{tcolorbox}

\section{Proofs}

\begin{lemma}[Exponential convergence]\label{lem:main}
Fix $t_0>0$ such that $\zeta(\sigma+it_0)\neq0$ for all $\sigma\in[0,1]$.
Set $m(t_0):=\min_{\sigma\in[0,1]}|\zeta(\sigma+it_0)|^2>0$.
Every solution of~\eqref{eq:flow} satisfies
$|\sigma(t)-\tfrac12|\le|\sigma_0-\tfrac12|\,e^{-m(t_0)\,t}$.
\end{lemma}

\begin{proof}
Let $u=\sigma-\tfrac12$.
Then $\dot u=-u|\zeta|^2$ and
$\tfrac{d}{dt}(u^2/2)=-u^2|\zeta|^2\le-m(t_0)\,u^2$.
By Gr\"onwall: $u(t)^2\le u(0)^2e^{-2m(t_0)t}$.
Taking square roots completes the proof.
(Part (i): $m(t_0)>0$ by the extreme value theorem and the
zero-free hypothesis.)
\end{proof}

\begin{proof}[Proof of Theorem~\ref{thm:kv}]
By Vinogradov--Korobov, $\zeta$ has no zeros for $\sigma>\sigma_{\mathrm{KV}}$.
Ford~\cite{ford} gives the sub-convexity bound
$|\zeta(\sigma+it_0)|^{-1}\le C_{\mathrm{KV}}(\log t_0)^{2/3}$ in this
region, so $m_{\mathrm{KV}}(t_0)>0$ and Lemma~\ref{lem:main} applies.
\end{proof}

\begin{remark}[Why no global lower bound exists]\label{rem:nobound}
A \emph{universal} $t_0$-independent lower bound
$m(t_0)\ge c>0$ cannot hold.  Non-trivial zeros cluster at
arbitrarily large imaginary parts; whenever a zero sits at height $t_0$
the minimum $m(t_0)=0$ by definition.  On nearby zero-free verticals
$m(t_0)$ can be driven arbitrarily close to zero by continuity.

What can exist—and what the flow argument requires—is a
\emph{window-dependent} bound: for every compact subset of the
zero-free region there exists $c(K)>0$ such that $m(t_0)\ge c(K)$
for all $t_0$ whose vertical segment stays inside $K$.  The present
paper produces such a bound unconditionally inside the Korobov--Vinogradov
strip (Theorem~\ref{thm:kv}).  Extending a positive window all the way
to the critical line $\sigma=\tfrac12$—for \emph{every} height—is
precisely the Riemann Hypothesis.
\end{remark}

\begin{proof}[Proof of Corollary~\ref{cor:equiv}]
($\Rightarrow$) If RH holds, no zeros exist off $\sigma=\tfrac12$;
on each compact segment $[0,1]\times\{t_0\}$ continuity of $\zeta$
and the absence of zeros gives $m(t_0)>0$ by the extreme value theorem.
($\Leftarrow$) If $m(t_0)>0$ for all $t_0$, Lemma~\ref{lem:main}
applies: all orbits converge to $\sigma=\tfrac12$.
The only fixed points of~\eqref{eq:flow} are Type~A ($\sigma=\tfrac12$)
and Type~B (zeros of $\zeta$); if $m(t_0)>0$ there are no Type~B fixed
points off $\sigma=\tfrac12$, so RH holds.
\end{proof}

\begin{remark}[Any improvement in zero-density tightens the bound]
The argument shows that any future lower bound
$m(t_0)\ge f(t_0)>0$ valid for $\sigma\in[0,\sigma_{\mathrm{KV}})$
automatically gives exponential convergence of the flow in that wider
strip, without changing the proof.
Zero-density tools (Huxley~\cite{huxley}, Bourgain--Gamburd--Sarnak,
Heath-Brown mean-value theorems) are the natural candidates.
\end{remark}

\section*{Appendix A: Numerical Evidence}

Over 60 sample starting points ($\sigma_0\in\{0.2,\ldots,0.8\}$,
first 10 non-trivial zeros), no off-critical fixed points were found.
Away from zero imaginary parts the measured halving time matches the
bound $\ln2/m_{\mathrm{KV}}$ to within 5\%; near zeros the flow slows
as expected.

\section*{Appendix B: Road-Map for Uniform Bounds Below $\sigma_{\mathrm{KV}}$}

The gap to an unconditional proof of RH is precisely a uniform lower
bound $m(t_0)>0$ for $\sigma\in[0,\sigma_{\mathrm{KV}})$.
Existing tools that could supply this (in increasing strength):

\begin{enumerate}
\item \textbf{Littlewood 1924, Ford 2002}: quantitative bounds in the
      zero-free strip give $m_{\mathrm{KV}}$ above; no extension below.
\item \textbf{Huxley's density estimates}: bounds on
      $N(\sigma,T):=\#\{\rho:\operatorname{Re}\rho>\sigma,\,|\operatorname{Im}\rho|<T\}$;
      sparse zeros $\Rightarrow$ $|\zeta|$ cannot be too small too often.
\item \textbf{Bourgain--Gamburd--Sarnak flattening}: $L^2$ mass-concentration
      bounds could prevent $|\zeta|^2$ approaching zero on a full set.
\item \textbf{Heath-Brown discrete mean-value theorems}: averaged lower
      bounds on $|\zeta(\sigma+it)|$ integrated over $t\in[T,2T]$.
\end{enumerate}

Each of these tools, if made pointwise rather than averaged, would
directly supply the uniform bound needed to close the flow argument.

\section*{Appendix C: Which Estimates Could Close the Inner Collar}

The open analytic question (eq.~\ref{eq:open}) asks for a lower bound
$|\zeta(\sigma+it_0)|\ge e^{-c(\log t_0)^{2/3}(\log\log t_0)^{1/3}}$
on every zero-free vertical segment.
We survey the tools in ascending order of strength.

\medskip\noindent\textbf{C.1 — Korobov--Vinogradov (current best, unconditional).}
For $\sigma\ge\sigma_{\mathrm{KV}}(t_0)=1-c(\log t_0)^{-2/3}(\log\log t_0)^{-1/3}$,
the zero-free region gives $|\zeta|^{-1}=O((\log t)^{2/3})$~\cite{ford},
so $m_{\mathrm{KV}}=\Omega((\log t)^{-4/3})>0$.
Covers $\sigma\in[\sigma_{\mathrm{KV}},1]\approx[0.999,1]$ at $t=10^6$.
\emph{Gap:} leaves $\sigma\in[\tfrac12,\sigma_{\mathrm{KV}})$ untouched.

\medskip\noindent\textbf{C.2 — Huxley's zero-density estimate.}
Huxley~\cite{huxley} gives $N(\sigma,T)=O(T^{A(1-\sigma)}(\log T)^B)$
with $A=12/5$.
A lower bound on $|\zeta|$ would follow if the zero count is so sparse
that $|\zeta|$ cannot approach zero everywhere on a vertical segment.
\emph{What is needed:} convert density to a pointwise lower bound.
Currently the density estimate is averaged over $T$, not pointwise in $t_0$.

\medskip\noindent\textbf{C.3 — Sub-convexity bounds (Weyl, Burgess, Bourgain).}
Upper bounds $|\zeta(\tfrac12+it)|\le t^\mu$ with $\mu<\tfrac14$ (Bourgain:
$\mu\le\tfrac{13}{84}$) constrain how large $|\zeta|$ can be,
but do not directly supply a lower bound.
\emph{Potential route:} if a zero existed at $\sigma_0+it_0$ with
$\sigma_0\ne\tfrac12$, the functional equation forces a paired zero at
$1-\sigma_0+it_0$; controlling $|\zeta|$ between the two might produce a
contradiction via a zero-free argument.

\medskip\noindent\textbf{C.4 — Mean-value theorems (Heath-Brown).}
Heath-Brown's discrete mean-value theorem bounds
$\sum_n|\zeta(\tfrac12+it_n)|^{2k}$.
A pointwise lower bound of the form
$\min_{|t-t_0|<1}|\zeta(\sigma+it)|^2\ge f(t_0)$
would follow from a lower bound on the mean combined with a Lipschitz
control on $\zeta'$.  \emph{Status:} not yet extracted from the literature.

\medskip\noindent\textbf{Summary.}
\begin{center}
\begin{tabular}{lll}
\hline
Tool & What it gives & What is still needed \\
\hline
Korobov--Vinogradov & $m>0$ for $\sigma\ge\sigma_{\mathrm{KV}}$ & Push to $\sigma>\tfrac12$ \\
Huxley density & Sparse zeros on average & Pointwise lower bound \\
Sub-convexity & Upper bound on $|\zeta|$ & Convert to lower bound \\
Heath-Brown mean & Averaged $|\zeta|^{2k}$ & Pointwise from average \\
\hline
\end{tabular}
\end{center}
The common obstacle is the gap between \emph{averaged} and
\emph{pointwise} estimates.
Any tool that produces a pointwise lower bound on $|\zeta(\sigma+it_0)|$
for individual $t_0$ (not averaged over $t_0$) would directly supply
the inner-collar bound and close the flow argument.

\begin{thebibliography}{9}
\bibitem{titchmarsh}
  E.~C. Titchmarsh (rev.\ D.~R. Heath-Brown),
  \textit{The Theory of the Riemann Zeta-Function}, 2nd~ed.,
  Oxford, 1986.
\bibitem{ford}
  K.~Ford,
  Zero-free regions for the Riemann zeta function,
  \textit{Number Theory for the Millennium II} (2002), 25--56.
\bibitem{huxley}
  M.~N. Huxley,
  On the difference between consecutive primes,
  \textit{Invent.\ Math.} \textbf{15} (1972), 164--170.
\bibitem{gronwall}
  T.~H. Gr\"onwall,
  Note on the derivatives with respect to a parameter \ldots,
  \textit{Ann.\ Math.} \textbf{20} (1919), 292--296.
\bibitem{tig}
  B.~Sanders, \textit{TIG Operator Specification (TSML)},
  Zenodo \href{https://doi.org/10.5281/zenodo.18852047}{10.5281/zenodo.18852047},
  2026.
\end{thebibliography}

\end{document}
